% GENERATED FILE. Edit canonical lessons, then run npm run generate:books.
% canonical-source-hash: fnv1a64:a7536ceac3748621
% canonical-lessons: HI-C130-hathaura, HI-C130-kil, HI-C130-kaici, HI-C130-dhaga, HI-C130-kanghi

\chapter{Hammer, Nail, Scissors, Thread, Comb}
\label{ch:hi-hammer-nail-scissors-thread-comb}
\hlchaptermodality{130}

\begin{chapteropening}
\textbf{By the end of this chapter:} \emph{I can say a hammer, a nail, scissors, thread and a comb.}

Complete the last lesson of chapter 130: i can say a hammer, a nail, scissors, thread and a comb.
\end{chapteropening}

\section[hathauṛā]{\hi{हथौड़ा} (hathauṛā) — a hammer}

\label{lesson:HI-C130-hathaura}



\begin{quote}
Before the new one: say the Hindi for a curtain, then the Hindi for a bucket.
\end{quote}

\subsection*{You'll want to know: \hi{हथौड़ा}}
\textbf{\hi{हथौड़ा}} — \emph{hathauṛā} — \textquotedblleft{}a hammer\textquotedblright{}.

\begin{grammarlens}[title={What you've built}]
One more word for this chapter.
\end{grammarlens}

\subsection*{Your turn}
\begin{itemize}
\raggedright
  \item \emph{Say it:} \emph{hathauṛā}
  \item \emph{Say it:} \emph{hathauṛā}, once more
  \item \emph{Say it:} say \emph{hathauṛā}
  \item \emph{Recall:} say the Hindi for a curtain, then the Hindi for a bucket, then say \emph{hathauṛā} again
\end{itemize}

\begin{tcolorbox}[breakable,colback=teal!4,colframe=teal!35!black,title={Before you move on}]
What does \hi{हथौड़ा} mean? (\textquotedblleft{}A hammer\textquotedblright{}.) Say it once more.
\end{tcolorbox}
\section[kīl]{\hi{कील} (kīl) — a nail (to hammer)}

\label{lesson:HI-C130-kil}



\begin{quote}
Before the new one: say the Hindi for a bucket, then the Hindi for a hammer.
\end{quote}

\subsection*{You'll want to know: \hi{कील}}
\textbf{\hi{कील}} — \emph{kīl} — \textquotedblleft{}a nail (to hammer)\textquotedblright{}.

\begin{grammarlens}[title={What you've built}]
One more word for this chapter.
\end{grammarlens}

\subsection*{Your turn}
\begin{itemize}
\raggedright
  \item \emph{Say it:} \emph{kīl}
  \item \emph{Say it:} \emph{kīl}, once more
  \item \emph{Say it:} say \emph{kīl}
  \item \emph{Recall:} say the Hindi for a bucket, then the Hindi for a hammer, then say \emph{kīl} again
\end{itemize}

\begin{tcolorbox}[breakable,colback=teal!4,colframe=teal!35!black,title={Before you move on}]
What does \hi{कील} mean? (\textquotedblleft{}A nail (to hammer)\textquotedblright{}.) Say it once more.
\end{tcolorbox}
\section[kaĩcī]{\hi{कैंची} (kaĩcī) — scissors}

\label{lesson:HI-C130-kaici}



\begin{quote}
Before the new one: say the Hindi for a hammer, then the Hindi for a nail.
\end{quote}

\subsection*{You'll want to know: \hi{कैंची}}
\textbf{\hi{कैंची}} — \emph{kaĩcī} — \textquotedblleft{}scissors\textquotedblright{}.

\begin{grammarlens}[title={What you've built}]
One more word for this chapter.
\end{grammarlens}

\subsection*{Your turn}
\begin{itemize}
\raggedright
  \item \emph{Say it:} \emph{kaĩcī}
  \item \emph{Say it:} \emph{kaĩcī}, once more
  \item \emph{Say it:} say \emph{kaĩcī}
  \item \emph{Recall:} say the Hindi for a hammer, then the Hindi for a nail, then say \emph{kaĩcī} again
\end{itemize}

\begin{tcolorbox}[breakable,colback=teal!4,colframe=teal!35!black,title={Before you move on}]
What does \hi{कैंची} mean? (\textquotedblleft{}Scissors\textquotedblright{}.) Say it once more.
\end{tcolorbox}
\section[dhāgā]{\hi{धागा} (dhāgā) — thread}

\label{lesson:HI-C130-dhaga}



\begin{quote}
Before the new one: say the Hindi for a nail, then the Hindi for scissors.
\end{quote}

\subsection*{You'll want to know: \hi{धागा}}
\textbf{\hi{धागा}} — \emph{dhāgā} — \textquotedblleft{}thread\textquotedblright{}.

\begin{grammarlens}[title={What you've built}]
One more word for this chapter.
\end{grammarlens}

\subsection*{Your turn}
\begin{itemize}
\raggedright
  \item \emph{Say it:} \emph{dhāgā}
  \item \emph{Say it:} \emph{dhāgā}, once more
  \item \emph{Say it:} say \emph{dhāgā}
  \item \emph{Recall:} say the Hindi for a nail, then the Hindi for scissors, then say \emph{dhāgā} again
\end{itemize}

\begin{tcolorbox}[breakable,colback=teal!4,colframe=teal!35!black,title={Before you move on}]
What does \hi{धागा} mean? (\textquotedblleft{}Thread\textquotedblright{}.) Say it once more.
\end{tcolorbox}
\section[kaṅghī]{\hi{कंघी} (kaṅghī) — a comb}

\label{lesson:HI-C130-kanghi}



\begin{quote}
Before the new one: say the Hindi for scissors, then the Hindi for thread.
\end{quote}

\subsection*{You'll want to know: \hi{कंघी}}
\textbf{\hi{कंघी}} — \emph{kaṅghī} — \textquotedblleft{}a comb\textquotedblright{}.

\begin{grammarlens}[title={What you've built}]
One more word for this chapter.
\end{grammarlens}

\subsection*{Your turn}
\begin{itemize}
\raggedright
  \item \emph{Say it:} \emph{kaṅghī}
  \item \emph{Say it:} \emph{kaṅghī}, once more
  \item \emph{Say it:} say \emph{kaṅghī}
  \item \emph{Recall:} say the Hindi for scissors, then the Hindi for thread, then say \emph{kaṅghī} again
\end{itemize}

\begin{tcolorbox}[breakable,colback=teal!4,colframe=teal!35!black,title={Before you move on}]
What does \hi{कंघी} mean? (\textquotedblleft{}A comb\textquotedblright{}.) Say it once more.
\end{tcolorbox}
